%!TEX ROOT = thesis.tex
\chapter{Conclusion}
\hspace{12.7mm}This study aims to develop cost-effective machine learning models for the early prediction of cognitive impairment using accessible data. The study presents a methodological framework designed to enhance the generalizability, explainability, and clinical utility of the developed models. The outcome of this study is designed to offer a cost-effective diagnosis tool that can be used reliably in primary care units where access to expensive and comprehensive diagnosis tools is limited. Ultimately, this research contributes to efforts to reduce the social, economic, and healthcare burden caused by the growing prevalence of cognitive impairment and dementia.

\section{Summary of Main Findings}
This study utilized a hierarchical classification approach to classify the impaired cases first before identifying the severity, cognitive impairment but not dementia (\acrshort{CIND}) or dementia, and the etiology of cognitive impairment, Alzheimer's disease (\acrshort{AD}), or other etiologies (non-AD). It also classifies current and future cognitive impairment separately. This approach improved classification accuracy and generalizability by addressing data complexity challenges such as class imbalance and class overlap. In addition, it improved the explainability and clinical utility of the developed models.

The proposed cost-effective models were developed using a methodological framework, which resulted in key findings supporting the objectives of this study. The proposed ensemble feature selection method consistently reduced the number of required features by approximately 82\% to 90\% across all classification tasks, retaining only 8 to 12 features from the original sets of 66 to 77 features. This resulted in reduced class overlap and computational cost while maintaining high diagnostic accuracy. The measures of instrumental activities of daily living were identified as the most important predictors across all cognitive status classification tasks, \hl{which is consistent with prediction models from different cohorts} \cite{Cui2024Prediction, Ren2025Using}.

Furthermore, the proposed overlap-guided resampling method significantly reduced data complexity. The proposed method improved the classification accuracy by balancing the recall of both targeted classes. For instance, in the future CIND vs. dementia prediction task, the recall of dementia in the CIND class improved from the baseline of 0.656 to 0.688. This resulted in an improvement in overall performance, with the geometric mean (\acrshort{G-mean}) increasing from 0.711 to 0.730. The proposed method outperformed standard data resampling methods such as the Synthetic Minority Over-sampling Technique (\acrshort{SMOTE}), which failed to improve performance and even led to worsening of the data complexity.

The novel model selection algorithm successfully identified hyperparameters that significantly increased training and prediction speeds compared to hyperparameters that produce the highest performance. The prediction throughput was accelerated up to 3 times with a minimal decrease in performance below 1\%. An evaluation of different machine learning algorithms showed that the logistic regression (\acrshort{LR}) model achieved superior computational efficiency and comparable predictive performance to complex, non-linear models such as random forest, support vector machines, and light gradient boosting machine.

The systematic data preprocessing methods and the comprehensive model selection resulted in cost-effective and generalizable LR models. The models maintained high predictive performance when tested on the independent Alzheimer's Disease Neuroimaging Initiative (\acrshort{ADNI}) dataset and unseen data from different research centers in the National Alzheimer's Coordinating Center (\acrshort{NACC}) Uniform Data Set (\acrshort{UDS}). The models achieved area under the curve (\acrshort{AUC}) scores of 0.825 and 0.841 to classify current cognitive status in the ADNI dataset as normal cognition (\acrshort{NC}) vs. cognitive impairment (CI) and CIND vs. dementia, respectively. To predict cognitive status up to 5 years in the future, the models scored 0.778 for the NC vs. CI task and 0.780 for the CIND vs. dementia task. However, the LR model that classifies the etiology of cognitive impairment (AD vs. non-AD) achieved modest performance, with an AUC score of 0.670. This finding confirms that accessible data alone is insufficient to differentiate specific causes of dementia. The models developed in this study outperformed those of similar studies in the literature that use accessible data to predict CI. \hl{They also performed within the range of the reported clinical accuracy of the underlying functional assessments, supporting the reliability of the obtained performance.}

Beyond classification accuracy, the Decision Curve Analysis (\acrshort{DCA}) demonstrated that using the proposed models provides a higher net benefit than default strategies (treat all and treat none) across a wide range of risk thresholds in the external datasets, ranging between 0.11 and 0.99. In addition, the LR models were perfectly calibrated for the NACC cross-center, with a calibration error below 0.05 for most of the classification tasks. However, the models showed significant miscalibration when applied to the ADNI dataset.

The explainability of the developed models at multiple levels revealed important findings. Local explanations revealed that patients with the same predicted risk probability often had completely different underlying risk factors. Moreover, the counterfactual model successfully generated valid counterfactual explanations for over 86\% of the cases in the external ADNI dataset. The counterfactual model was able to generate counterfactual explanations for all cases, identifying modifiable risk factors to prevent future risk of CI. The counterfactual analysis showed that for future risk prediction, reversing the risk often required modifying only 1 to 3 features, whereas reversing a current diagnosis required significantly more changes. This finding provides data-driven evidence that early prevention of CI risk is much more feasible and effective than reversing established impairment.

\section{Significance and Contribution of the Study}
Based on these findings, significant contributions to both theory and practice were identified. The results showed how effective the proposed methodological framework is in developing cost-effective, generalizable, and reliable models for the early prediction of cognitive impairment. The proposed feature selection and data resampling methods demonstrated that the development of cost-effective models should be data-driven. This becomes especially critical when building the models using features with low discriminative power, such as basic demographic and functional data. The results indicated that when data complexity is reduced, simple models can accurately classify target classes as effectively as complex and ensemble models. Therefore, more efforts should be directed towards refining data preprocessing techniques compared to building sophisticated classification models.

Furthermore, for the CI prediction models to be used in diverse clinical settings with varying resource levels, the model fine-tuning and selection process should involve evaluation of the model's complexity and computational cost instead of mere assessment of classification performance. The model selection method in this study allowed the selection of models that are efficient and accurate. This reduced the development and operational cost of the models and promoted their deployment to different real-world application scenarios, including remote healthcare facilities with limited computing resources.

The application of systematic data preprocessing and comprehensive model selection methods could lead to the development of generalizable models. This is supported by the models' performance in the internal cross-validation and external validations. \hl{It suggests that cost-effective models may reliably predict cognitive status across similar research centers and study protocols.} Furthermore, the analysis of the results in this study concluded that while accessible data are sufficient for the development of accurate models to classify the severity of CI, they are limited in classifying the etiology of CI. This finding highlights the need for novel cost-effective indicators that facilitate the accurate classification of CI subtypes beyond the use of expensive and invasive data.

Moreover, the developed models are able to produce estimates that accurately reflect the actual risk of CI and dementia across a broad range of probability thresholds. This makes the developed models clinically reliable and useful for various use cases. The models can be used in use cases that require accurate prediction of CI and dementia at low-risk thresholds. This applies to population-level screening where sensitive risk prediction is required to identify potential cases, such as routine health check-ups and community health programs. On the other hand, the models can also be used for the prediction of targeted risk in specialized memory clinics where precision is essential.

The multilayered model explainability method in this study also supports the clinical utility and reliability of the models. The counterfactual explanation demonstrated the great potential of CI and dementia risk prevention through actionable and insightful explanation of the model prediction. Overall, this study encourages the development of interpretable models and the implementation of personalized explanations, including counterfactual explanations, that go beyond the global explanation of the model prediction. The study suggests that explanation of individual model predictions is crucial to support personalized recommendations and communication of risk factors, which improves risk management and reliability between clinicians and patients.

\hl{The developed model can serve as an early screening tool to automatically estimate the potential risk of CI and dementia. The model estimation can be generated during a routine primary care visit using already collected data, adding only a minor step at minimal extra cost. Based on the model estimation, clinicians and healthcare professionals can decide whether further specialized and expensive assessments, such as cognitive tests or neuroimaging, are required. The main added cost would be the follow-up of the false positive cases, which could be acceptable for a screening tool. Additionally, clinicians can design tailored intervention plans based on local and counterfactual explanations of the model predictions.}

The proposed LR models are efficient and require low computing resources to operate. Therefore, the models can be deployed as a service on cloud-based platforms with minimal operational cost, or they can be hosted on local servers or even edge devices for areas with limited internet access. In addition, primary care units can integrate the models with their healthcare electronic systems and database management systems to generate real-time predictions based on updated patient data.

\newpage
This integration optimizes the diagnosis workflow by reducing the time and resources required for preliminary assessments. Moreover, it significantly reduces the cost and removes the barriers to accessing early prediction of CI risk, especially in low- and middle-income countries (\acrshort{LMICs}), where access to specialized tools and professional resources is limited. Furthermore, the ease of integration enables widespread access to the models through telemedicine platforms. This enables healthcare providers to provide an early assessment via video calls, phone calls, or online platforms, thereby reducing the need for in-person visits. Moreover, the models can be embedded into mobile applications and run on personal edge devices, empowering individuals to personally monitor their cognitive status in real time based on data retrieved from electronic databases.

In summary, using the developed cost-effective model as an early screening tool not only improves the efficiency of the diagnostic process but also simplifies it, making it directly accessible to primary care providers, individuals, and caregivers. Therefore, the outcome of this study aims to reduce the increasing burden of CI and dementia by offering access to early risk prediction and prevention. In addition, the study contributes towards the United Nations Sustainable Development Goals (\acrshort{SDGs}) by promoting good health and well-being (SDG 3) and offering equal accessibility to early prediction and prevention of cognitive impairment (SDG 10).

\section{Limitations of the Study}
Despite the significant contributions of this study, there are several limitations that need to be acknowledged. First of all, the datasets used in this study present critical constraints that might impact the generalizability of the obtained results. NACC UDS and ADNI are longitudinal clinical measurements and diagnostic data. They are collected from participants in the United States who are mostly white, highly educated, and actively looking for medical advice. Although these datasets were used because they offer rich data for this research, the findings may not be generalizable to the general public or to more diverse global populations, including LMICs where education levels and healthcare access are different. Therefore, the models may underperform if applied to populations with different prevalence of CI and socio-economic backgrounds.

This study utilized data collected through different versions of NACC UDS. The evolution of the UDS forms introduced systematic missing data and a change in the definition and criteria for mild cognitive impairment and AD. This led to the exclusion of potentially useful variables that were not collected in all versions and might introduce inconsistencies in the target labels, potentially affecting the model's stability. In addition, ADNI was designed to specifically identify biomarkers for AD; thus, it excludes other conditions, such as vascular dementia. This caused the external validation of the CI subtype classification model to be limited to testing only the generalizability of the model across NACC research centers.

The proposed data preprocessing and model selection methods possess specific limitations that might limit their broad adoption. In general, the application of these methods may increase the computational cost and training time compared to other standard methods. With respect to feature selection, the maximum Fisher's Discriminant ratio (\acrshort{F1}) is biased towards selecting features that improve linear separability and reduce class overlap. Using F1 as a threshold in the proposed feature selection method could exclude non-linear features that increase predictive performance. Although this method proved to enhance the model performance in this study, its generalizability to other datasets must be tested.

Furthermore, the diagnosis of cognitive impairment in the NACC study is made based on clinical assessments rather than modifiable lifestyle factors, such as diet, sleep, and physical activity. This could influence the selected features, which were based primarily on functional assessments and dementia clinical ratings. Consequently, counterfactual explanations could become less actionable. For instance, it is much harder to reverse a deficit in managing finances than it is to improve a sleep habit.

Regarding the proposed data resampling method, the use of Edited Nearest Neighbor (\acrshort{ENN}) might have excluded useful boundary instances from both target classes that define clear decision boundaries. By removing these difficult samples from the training data, the model may be ineffective in diagnosing complex cases in real clinical settings. In terms of the model selection method, the efficiency of the models was mainly measured based on computational throughput (samples per second). This measure does not provide a holistic assessment of the efficiency of the models', and it is highly dependent on the computational resources used to train and infer the models. The efficiency assessment should also include the footprint of the memory and the size of the model.

This study is designed based on the hierarchical classification approach, which might have caused error propagation. For instance, a dementia patient who is wrongly classified as normal by the first-level model is not included in the second-level model for further evaluation. This means that patients are only classified as CIND or dementia if they are accurately classified as impaired. This might also lead to complex and conflicting explanations. Although the hierarchical classification approach simplifies the data complexity and reflects the clinical workflow, it possibly introduces complexity to the decision process. Lastly, the performance of the model explainability methods was evaluated quantitatively and qualitatively in this study. \hl{However, clinical validation is still required to evaluate the accuracy of the model predictions and the effectiveness of the generated explanations in a real diagnostic workflow.}

\section{Recommendations for Future Work}
While this research successfully demonstrated the feasibility of cost-effective prediction of CI, the following recommendations are suggested for further improvements. Firstly, it is recommended to evaluate the proposed methodological framework on a community-based dataset from LMICs. This work is crucial in assessing cost-effective models in low-resource environments where they are most needed. Furthermore, since the current study conducted a retrospective analysis of existing datasets, \hl{it is advisable to test the developed models in real-time primary care settings to clinically validate their performance in the early prediction and prevention of the risk of CI.}

The proposed feature selection and data resampling method are tested only on the NACC UDS and the context of this study. Consequently, the methods must be tested on various datasets and various classification tasks to evaluate their generalizability and identify areas of improvement. Enhancements might include the use of multiple class overlap measures for feature selection and the incorporation of more effective data resampling methods. \hl{Additionally, clinical evaluation should be integrated as part of the feature selection process so that the clinical applicability of the ranked features is ensured.} Furthermore, future research should consider incorporating different accessible data modalities that can improve the accuracy of cost-effective models without relying on expensive imaging. For instance, further efforts could incorporate data based on speech and handwriting analysis, which can be collected non-invasively using smartphones.

The current approach uses baseline data to predict CI at a fixed future point (5 years in the future). For further improvement, it is suggested to explore the modeling of cognitive status progression to predict the time of conversion from normal cognition to CI or dementia. This could offer clinicians valuable insights to support more precise risk prevention plans. Moreover, it is beneficial to utilize longitudinal input data to capture the temporal patterns of cognitive decline, providing a more accurate representation of disease progression.

Finally, generative artificial intelligence, specifically small language models (SLMs), could be utilized to enhance the explanation of the predicted risk of CI. SLMs can be used to generate detailed and human-understandable explanations that can facilitate easier communication between clinicians and patients.

At the end of this research work, it is hoped that the framework proposed in this study will serve as a foundation for scalable, accessible, and equitable cognitive impairment screening tools globally.